\begin{abstract}
Recent methods for collision avoidance in Vision-Language-Action (VLA) models report high Collision Avoidance Rate (CAR) on safety benchmarks and claim effective obstacle avoidance. We show that these metrics can be systematically contaminated by action collapse: the model ceases to produce meaningful motion, and the resulting inaction is counted as successful avoidance. On the SafeLIBERO benchmark, a Signed Distance Field (SDF)-augmented VLA reports 91.5\% CAR, yet 99.6\% of its safe episodes exhibit end-effector displacement below 0.1\,mm. This contamination is not method-specific; a baseline without any collision objective achieves 82\% CAR with 99.4\% of safe episodes below the same threshold, and a more capable model shows the same pattern. We propose displacement-filtered CAR (CAR$_\text{f}$), which excludes sub-threshold episodes, as a corrective metric. Further analysis reveals two structural issues. First, Task Success Rate (TSR) and CAR are anti-correlated: models that move the robot through the workspace collide, while collapsed models appear safe. Second, ground-truth SDF oracle experiments demonstrate that end-effector-only refinement cannot prevent full-body collisions, as SDF gradients that repel the end-effector simultaneously drive arm links into obstacle volumes. Preliminary results indicate that training-time SDF supervision on a sufficiently capable VLA can produce genuine collision avoidance with zero frozen-action rate. Our findings establish that displacement filtering is necessary for meaningful VLA safety evaluation.
\end{abstract}
